% !TEX program = pdflatex
% !BIB program = bibtex
\documentclass[11pt]{article}

% Change "review" to "final" to generate the final (sometimes called camera-ready) version.
% Change to "preprint" to generate a non-anonymous version with page numbers.
\usepackage[review]{acl}

% Standard package includes
\usepackage{times}
\usepackage{latexsym}

% For proper rendering and hyphenation of words containing Latin characters (including in bib files)
\usepackage[T1]{fontenc}
% For Vietnamese characters
% \usepackage[T5]{fontenc}
% See https://www.latex-project.org/help/documentation/encguide.pdf for other character sets

% This assumes your files are encoded as UTF8
\usepackage[utf8]{inputenc}

% This is not strictly necessary, and may be commented out,
% but it will improve the layout of the manuscript,
% and will typically save some space.
\usepackage{microtype}

% This is also not strictly necessary, and may be commented out.
% However, it will improve the aesthetics of text in
% the typewriter font.
\usepackage{inconsolata}

%Including images in your LaTeX document requires adding
%additional package(s)
\usepackage{graphicx}
\usepackage{amsmath}
\usepackage{amssymb}
\usepackage{booktabs}
\usepackage{multirow}
\usepackage{float}
\usepackage{placeins}
\usepackage{stfloats}
\usepackage{adjustbox}
\usepackage[nameinlink,noabbrev]{cleveref}
\emergencystretch=2em
\hfuzz=20pt
\vfuzz=20pt
\hbadness=10000
\vbadness=10000
% If the title and author information does not fit in the area allocated, uncomment the following
%
%\setlength\titlebox{<dim>}
%
% and set <dim> to something 5cm or larger.

\usepackage{soul}
\usepackage{color, xcolor}
\definecolor{lightblue}{RGB}{0, 191, 255}
\sethlcolor{lightblue}
\soulregister\cite7
\soulregister\eqref7
\soulregister\ref7 
\newcommand{\mathcolorbox}[2]{\colorbox{#1}{$\displaystyle #2$}}
\newcommand{\clr}[1]{\textcolor{blue}{#1}}

\title{A Multi-cluster Boundary Learning Method for Out-of-scope Intent Detection via laernable MiniLM Embedding}

% \author{
%   \textbf{Yihong Xu\textsuperscript{1}},
%   \textbf{Mingyu Kang\textsuperscript{1*}},
%   \textbf{Linyuan L\"u\textsuperscript{1*}}\\
%   \textsuperscript{1}University of Science and Technology of China, Hefei, China\\
%   \texttt{xuyihong@mail.ustc.edu.cn}\\
%   \texttt{kangmingyu@ustc.edu.cn}\\
%   \texttt{linyuan.lv@ustc.edu.cn}
% }
% Author information can be set in various styles:
% For several authors from the same institution:
% \author{Author 1 \and ... \and Author n \\
%         Address line \\ ... \\ Address line}
% if the names do not fit well on one line use
%         Author 1 \\ {\bf Author 2} \\ ... \\ {\bf Author n} \\
% For authors from different institutions:
% \author{Author 1 \\ Address line \\  ... \\ Address line
%         \And  ... \And
%         Author n \\ Address line \\ ... \\ Address line}
% To start a separate ``row'' of authors use \AND, as in
% \author{Author 1 \\ Address line \\  ... \\ Address line
%         \AND
%         Author 2 \\ Address line \\ ... \\ Address line \And
%         Author 3 \\ Address line \\ ... \\ Address line}

% \author{First Author \\
%   Affiliation / Address line 1 \\
%   Affiliation / Address line 2 \\
%   Affiliation / Address line 3 \\
%   \texttt{email@domain} \\\And
%   Second Author \\
%   Affiliation / Address line 1 \\
%   Affiliation / Address line 2 \\
%   Affiliation / Address line 3 \\
%   \texttt{email@domain} \\}



%\author{
%  \textbf{First Author\textsuperscript{1}},
%  \textbf{Second Author\textsuperscript{1,2}},
%  \textbf{Third T. Author\textsuperscript{1}},
%  \textbf{Fourth Author\textsuperscript{1}},
%\\
%  \textbf{Fifth Author\textsuperscript{1,2}},
%  \textbf{Sixth Author\textsuperscript{1}},
%  \textbf{Seventh Author\textsuperscript{1}},
%  \textbf{Eighth Author \textsuperscript{1,2,3,4}},
%\\
%  \textbf{Ninth Author\textsuperscript{1}},
%  \textbf{Tenth Author\textsuperscript{1}},
%  \textbf{Eleventh E. Author\textsuperscript{1,2,3,4,5}},
%  \textbf{Twelfth Author\textsuperscript{1}},
%\\
%  \textbf{Thirteenth Author\textsuperscript{3}},
%  \textbf{Fourteenth F. Author\textsuperscript{2,4}},
%  \textbf{Fifteenth Author\textsuperscript{1}},
%  \textbf{Sixteenth Author\textsuperscript{1}},
%\\
%  \textbf{Seventeenth S. Author\textsuperscript{4,5}},
%  \textbf{Eighteenth Author\textsuperscript{3,4}},
%  \textbf{Nineteenth N. Author\textsuperscript{2,5}},
%  \textbf{Twentieth Author\textsuperscript{1}}
%\\
%\\
%  \textsuperscript{1}Affiliation 1,
%  \textsuperscript{2}Affiliation 2,
%  \textsuperscript{3}Affiliation 3,
%  \textsuperscript{4}Affiliation 4,
%  \textsuperscript{5}Affiliation 5
%\\
%  \small{
%    \textbf{Correspondence:} \href{mailto:email@domain}{email@domain}
%  }
%}

\begin{document}
\maketitle

\begin{abstract}
Intent detection is a critical task that bridges human intents and system actions in human-machine interaction systems. However, there still exist challenges for detecting out-of-scope (OOS) intents. (i) The traditional methods view the OOS intent detection as a multi-class classification, then the detection accuracy decreases as the class number of the known intents increases; (ii) LLM-embedding methods require large parameters, that makes them difficult to train and practically deploy. Thus, this work proposes a multi-cluster boundary learning method to detect OOS intents via MiniLM embedding (i.e., \texttt{all-MiniLM-L6-v2}) in an one-class classification workflow. The method learns the boundaries of multi-cluster embeddings generated by MiniLM from the training utterances, and then rejects the out-of-domain utterances as OOS intents. Experiments are conducted on public CLINC150, StackOverflow and Banking77 datasets. The results show that the method achieves the state-of-the-art OOS intent detection performance compared the other baselines. Ablation studies are also conducted and the results show that the used MiniLM can better adapt to the workflow and utterance embedding requirements. The code is available at supplementary materials.
\end{abstract}

\section{Introduction}
Intent detection is a critical yet challenging task for human-machine interaction systems. It maps user utterances to system actions, which bridges human intents and system behaviors~\citep{tur2010what,louvan2020recent,wolflein2025agents}. However, human intents are complex and cannot be fully enumerated in a closed-set intent detection system. That means if the system receives an out-of-scope (OOS) intent utterance but fails to reject it, that will trigger incorrect actions and induce harmful consequences, as shown in figure~\ref{fig:motivation}. Thus, a more important and difficult task is to detect the OOS intents and reject them in a timely manner, which is actually an open intent detection task~\citep{hoffman2024inferring, muzahid2024survey}.

\begin{figure}[t]
    \centering
    \includegraphics[width=\columnwidth]{figures/bg 02.png}
    \caption{A diagram of OOS intent detection}
    \label{fig:motivation}  
\end{figure}

There are three types of intent detection methods, including statistical methods, deep-learning methods and LLM-embedding methods. Among them, the statistical methods mainly use lexical or grammar-based features with statistical classifiers, such as SVM~\citep{haffner2003optimizing} and n-gram classifiers~\citep{tur2010what}. But these methods still depend on manually designed or task-dependent features. Then, the deep-learning methods learn semantic representations directly from utterances, e.g., RNN~\citep{liu2016attention} and pretrained BERT~\citep{chen2019bert}. However, these methods are mainly trained under the closed-set setting and cannot explicitly reject OOS intents. Thus, several open intent detection methods introduce explicit rejection mechanisms, including confidence and distance scoring~\citep{hendrycks2017baseline,lee2018simple}, contrastive representation learning~\citep{zeng2021modeling}, adaptive decision boundaries~\citep{zhang2021adaptive,zhang2023learning} and multi-granularity boundaries~\citep{li2025multi}. After that, LLM-embedding methods further use sentence embeddings and LLMs for intent detection and OOS rejection~\citep{zhang2024new,wang2024beyond}. Moreover, uncertainty-based routing combines lightweight models with LLMs to reduce unnecessary LLM inference~\citep{arora2024intent,zaera2025efficient}.

However, the existing methods still face two limitations. First, the OOS queries with similar semantic meaning as known intents are easily absorbed into the decision regions of the known intents~\citep{li2025multi}, but in fact they are OOS. Thus, if mix OOS labels into known intent labels, and simply view the OOS intent detection as a multi-class classification, the accuracy of OOS rejection would decrease as the class number increases. Second, although LLM-embedding methods improve the ability of semantic understanding, they still require very large parameters, that induces remarkable computational costs and high sensitivity of prompt designs~\citep{zaera2025efficient}. Thus, they are difficult to train and deploy in real-time resource-constrained dialogue systems.

To address these limitations, this work proposes a multi-cluster boundary learning method for OOS intent detection with a cascade workflow. The workflow decouples the complex open intent detection problem into multiple simpler stages. This allows the workflow to use a lightweight MiniLM instead of LLM wih large-scale parameters. Moreover, it reduces the OOS intent detection to a one-class classification, instead of a complex multi-class classification. Moreover, we found that MiniLM can embed the observation data into multiple clusters. Thus, if the boundaries of clusters are learned, the out-of-domain utterances can be identified as OOS intents.

Thus, the main contributions are as follows:
\begin{enumerate}
    \item[1.] A multi-cluster boundary learning method is proposed to realize OOS detection through learnable MiniLM embeddings. The learnable MiniLM adapts the representation space to the known intents, so that samples of the same intent become more compact while different intents remain separated. This provides a better representation space for boundary learning.
    
    \item[2.] A cascade workflow is proposed to first conduct OOS intent detection as an one-class classification, and then conduct known intent detection as a close-set multi-class classification. This means the OOS intent detection task can be addressed independently. 

    \item[3] Experiments are conducted on public real-world datasets for OOS intent detection task. The results show that the method achieves the state-of-the-art performance for OOS intent detection. Moreover, ablation studies are also conducted, and the results show that the used MiniLM, \texttt{all-MiniLM-L6-v2}, can better adapt to the workflow and utterance embedding requirements.  
\end{enumerate}


\section{Related Work}
The open intent detection methods can be mainly categorized into three types, i.e., statistical methods, deep-learning methods and LLM-embedding methods, according to the problem settings and the types of classifiers.   

\subsection{Statistical Methods for Intent Detection}
The statistical methods mainly relied on lexical patterns and statistical classifiers to detect intent classes. Among them, \citet{gorin1997how} detect call types using task-relevant grammar fragments learned from spoken utterances. Then, \citet{wang2002combination} combine statistical classifiers with a semantic grammar, but the grammar remains task-dependent. \citet{haffner2003optimizing} further use grammar fragments as SVM features and combine binary classifiers for multi-class classification. However, these methods still depend on manually designed features. Moreover, \citet{tur2010what} use word n-gram features with a discriminative classifier for intent determination, but their error analysis shows that local n-gram context remains limited for more complex utterances.

\subsection{Deep-Learning Methods for Intent Detection}
Then, the deep-learning methods learn semantic representations directly from utterances instead of relying on predefined lexical features. \citet{sarikaya2011deep} introduce Deep Belief Networks to learn hierarchical representations from unlabeled data and fine-tune them for call routing. \citet{liu2016attention} further use RNNs for joint intent detection and slot filling. Moreover, \citet{chen2019bert} introduce BERT for intent detection with pretrained contextual representations. However, these methods still assume a predefined intent set and cannot explicitly reject OOS inputs~\citep{hendrycks2017baseline, larson2019evaluation}.

Thus, several methods introduce explicit rejection mechanisms for unknown inputs. \citet{hendrycks2017baseline} use maximum softmax probability as a confidence score for OOD detection, and \citet{lee2018simple} further use Mahalanobis distance based on class-conditional Gaussian features. However, these methods mainly rely on confidence scores or predefined feature distributions. Then, \citet{shu2017doc} propose Deep Open Classification (DOC), which uses class-wise sigmoid classifiers and rejection thresholds for open text classification. \citet{zeng2021modeling} further use supervised contrastive learning to reduce intra-class variance and enlarge inter-class variance. But these methods do not explicitly learn class-specific geometric boundaries. Thus, \citet{zhang2021adaptive} propose Adaptive Decision Boundary (ADB), which learns an adaptive spherical boundary for each known class. \citet{zhang2023learning} further propose Distance-Aware Adaptive Decision Boundary (DA-ADB), which learns distance-aware intent representations before boundary optimization. However, these methods still model each known intent with a single spherical region. To model more complex intent distributions, \citet{li2025multi} propose Multi-granularity Open Intent Classification via Adaptive Granular-Ball Decision Boundary (MOGB), which uses adaptive granular-ball clustering and multiple local boundaries for each intent.

\subsection{LLM-Embedding Methods for Intent Detection}
Then, LLM-embedding methods are proposed to learn semantic representations for intent detection and OOS rejection. \citet{zhang2024new} fine-tune sentence transformers with an embedding reconstruction loss to reduce the dispersion of in-scope representations and improve OOS rejection. However, the rejection performance still depends on the learned embedding space. Thus, \citet{wang2024beyond} explore LLMs for OOS intent detection and show strong zero-shot and few-shot performance, but they remain less competitive than fully fine-tuned models under the full-resource setting. To reduce the inference cost, \citet{arora2024intent} combine sentence transformers with LLMs through uncertainty-based routing. Similarly, \citet{zaera2025efficient} invoke a fine-tuned LLM only for samples with high classifier uncertainty. However, these methods still require additional LLM inference for uncertain samples.

Then our previous work~\citep{xu2026multi} proposes a lightweight multi-cluster boundary method with MiniLM embeddings for OOS rejection. However, the MiniLM encoder is fixed, so its representation cannot be directly adapted for OOS detection.

Thus, the existing methods still face two problems. (i) the OOS intent detection is usually viewed as a multi-class classification and hardly separated from the known intent detection. Many existing methods perform OOS rejection and known intent discrimination in the same learned representation space. (ii) The multi-agent LLMs require remarkably large computation cost. And due to the reason (i), the current routing mechanism is not accurate. 

\section{Method}

\subsection{Cascade Workflow for Intent Detection}
\label{sec:cascade}

The proposed method uses a three-stage cascade, i.e., Gate $\rightarrow$ Router $\rightarrow$ Expert, as shown in Figure~\ref{fig:model_architecture}. Given an input utterance $x$, the Gate first determines whether $x$ belongs to the Known intent space. If $x$ is rejected, the system directly outputs an OOS label. Otherwise, the Router predicts its coarse-grained domain, and the corresponding Expert further predicts the fine-grained Known intent.

The three stages are defined as
\begin{align}
z &= G(x), && z\in\{\mathrm{Known},\mathrm{OOS}\}, \label{eq:gate_workflow}\\
d &= R(x), && d\in\mathcal{D}, \label{eq:router_workflow}\\
y &= E_d(x), && y\in\mathcal{Y}_d\subseteq\mathcal{Y}_{\mathrm ID},
\label{eq:expert_workflow}
\end{align}
where $\mathcal{D}$ is the domain set, $\mathcal{Y}_{\mathrm ID}$ is the Known intent set, and $\mathcal{Y}_d$ is the intent subset of domain $d$. The Router and Expert are executed only when $z=\mathrm{Known}$.

Thus, the Gate performs OOS detection before the Router and Experts perform known intent classification. Since the downstream task is closed-set, the Router and Experts are trained only with known intent samples using \texttt{SmolLM-135M} with LoRA~\citep{hu2022lora}. As a result, the Gate only needs to determine whether an input belongs to the Known intent space. Thus, a lightweight encoder is sufficient.

\begin{figure*}[t]
    \centering
    \includegraphics[width=\textwidth]{figures/overview 01.eps}
    \caption{Overview of the proposed Gate--Router--Expert workflow.}
    \label{fig:model_architecture}
\end{figure*}


\subsection{MiniLM-Based Representation Learning}
\label{sec:minilm_rep}

Since the Gate only needs to distinguish OOS from known intent samples, a lightweight MiniLM encoder, \texttt{all-MiniLM-L6-v2}, is used to construct the semantic representation space. Given an utterance $x$, MiniLM produces
\begin{equation}
e_0=\operatorname{MeanPool}\left(f_{\theta}(x)\right),
\qquad e_0\in\mathbb{R}^{h},
\label{eq:embed}
\end{equation}
where $f_{\theta}(\cdot)$ denotes the MiniLM encoder and $h=384$ is the embedding dimension.

A frozen MiniLM provides general sentence representations, but its embedding space is not directly optimized for OOS detection. Thus, we further adapt MiniLM using only known intent training samples. Selected Transformer blocks are updated, and a residual projection is applied:
\begin{equation}
e=e_0+
W_2\operatorname{GELU}
\left(W_1\operatorname{LN}(e_0)\right),
\label{eq:projection}
\end{equation}
where $\operatorname{LN}(\cdot)$ denotes layer normalization, and $W_1$ and $W_2$ are projection matrices. When the projection is disabled, $e=e_0$. Then, the embedding is normalized as $\bar e=e/\|e\|_2$.

To learn the known intent geometry, let $(x_i,y_i)$ denote the $i$-th training sample and $\mathcal{Y}_{\mathrm K}$ the Known intent set. For each intent $y$, its prototype is
\begin{equation}
p_y=
\frac{
\sum_{i:y_i=y}\bar e_i
}{
\left\|
\sum_{i:y_i=y}\bar e_i
\right\|_2
},
\label{eq:training_prototype}
\end{equation}
and $\delta_{i,y}=\|\bar e_i-p_y\|_2^2$ denotes the squared distance from sample $i$ to prototype $p_y$.

The representation objective is:
\begin{equation}
\mathcal{L}_{\mathrm{rep}}
=
\mathcal{L}_{\mathrm{ce}}
+\alpha\mathcal{L}_{\mathrm{intra}}
+\beta\mathcal{L}_{\mathrm{inter}},
\label{eq:rep_loss}
\end{equation}
\begin{align}
\mathcal{L}_{\mathrm{ce}}
&=
-\frac{1}{N}\sum_i
\log
\frac{\exp(-\delta_{i,y_i}/T)}
{\sum_{y\in\mathcal{Y}_{\mathrm K}}\exp(-\delta_{i,y}/T)},\\
\mathcal{L}_{\mathrm{intra}}
&=\frac{1}{N}\sum_i\delta_{i,y_i},\\
\mathcal{L}_{\mathrm{inter}}
&=
\frac{1}{N}\sum_i
\left[
\delta_{i,y_i}-\min_{y\neq y_i}\delta_{i,y}+m
\right]_{+}.
\label{eq:rep_terms}
\end{align}
where $N$ is the number of training samples, $\alpha$ and $\beta$ are loss weights, $T$ is the temperature, $m$ is the inter-intent margin, and $[\cdot]_+=\max(0,\cdot)$. Among them, $\mathcal{L}_{\mathrm{ce}}$ preserves known intent discrimination, $\mathcal{L}_{\mathrm{intra}}$ reduces within-intent distances, and $\mathcal{L}_{\mathrm{inter}}$ separates competing intents.

Thus, the adapted MiniLM produces compact representations for the same Known intent while maintaining inter-intent separation. No OOS samples are used for representation learning. The learned representation is then used for boundary learning.

\subsection{Boundary Learning for OOS Intent Detection}
\label{sec:boundary}

Based on the learned representations, the Gate constructs local acceptance regions for Known intents. For each intent $y$, one or multiple centers are used to represent its embedding distribution:
\begin{equation}
\mathcal{C}_y
=
\{c_{y,1},c_{y,2},\ldots,c_{y,K}\},
\label{eq:center_set}
\end{equation}
where $K$ is the number of centers. The $k$-th center is computed as
\begin{equation}
c_{y,k}
=
\frac{1}{|\mathcal{I}_{y,k}|}
\sum_{i\in\mathcal{I}_{y,k}}
\bar e_i,
\label{eq:center}
\end{equation}
where $\mathcal{I}_{y,k}$ denotes the samples assigned to the $k$-th center.

Then, the distance between an input embedding $\bar e$ and each center is calculated. Besides Euclidean distance, we use diagonal Mahalanobis distance to account for different variations across embedding dimensions:
\begin{equation}
d^{\mathrm M}_{y,k}(\bar e)
=
\sqrt{
\sum_{j=1}^{h}
\frac{(\bar e_j-c_{y,k,j})^2}
{v_{y,k,j}+\epsilon}
},
\label{eq:mahalanobis}
\end{equation}
where $v_{y,k,j}$ is the variance of the $j$-th dimension and $\epsilon$ is a small constant.

For each center, a local radius is estimated from the training distances:
\begin{equation}
r_{y,k}
=
\mu_{y,k}+\lambda\sigma_{y,k},
\label{eq:radius_mean_std}
\end{equation}
where $\mu_{y,k}$ and $\sigma_{y,k}$ are the mean and standard deviation of the distances, and $\lambda$ controls the boundary width. A quantile-based radius can also be used.

The normalized boundary score is defined as
\begin{equation}
b(\bar e)
=
\min_{y,k}
\frac{d_{y,k}(\bar e)}
{r_{y,k}},
\label{eq:normalized_union}
\end{equation}
where a smaller score indicates that the sample is closer to a Known-intent region. When ambiguity information is used, the final score becomes
\begin{equation}
s(\bar e)=b(\bar e)+\gamma a(\bar e),
\label{eq:fusion_score}
\end{equation}
where $a(\bar e)$ measures the ambiguity between the nearest intents and $\gamma$ controls its contribution. When $\gamma=0$, only the boundary score is used.

Finally, the Gate predicts
\begin{equation}
G(x)=
\begin{cases}
\mathrm{Known}, & s(\bar e)\leq\tau,\\
\mathrm{OOS}, & s(\bar e)>\tau,
\end{cases}
\label{eq:gate_decision}
\end{equation}
where $\tau$ is the decision threshold. Known samples are passed to the Router and Expert, while the remaining samples are rejected as OOS.

All boundary settings, including $K$, the distance function, radius estimator, score fusion, and threshold, are selected using only Known training and validation samples. No real or pseudo OOS samples or test samples are used for model selection.
\section{Experiments}
\subsection{Datasets}
Three datasets are used for performance evaluation, as shown in Table~\ref{tab:dataset_stats}. For each dataset, the samples from the known intents are split into training, validation, and test sets with a ratio of 6:1:3.


\textbf{CLINC150} is a large-scale multi-domain dataset with 150 in-scope intents and explicit OOS samples~\citep{larson2019evaluation}. The data source is available at \url{https://github.com/clinc/oos-eval}.

\textbf{StackOverflow} is a technical-domain short-text dataset with 20 intent classes and substantial semantic overlap across classes. The data source is available at \url{https://github.com/jacoxu/StackOverflow}.

\textbf{Banking77} is a fine-grained, single-domain dataset with 77 intent classes. It focuses on high-granularity classification in the financial domain. The data source is available at \url{https://github.com/PolyAI-LDN/task-specific-datasets}.

\begin{table}[H]
\centering
\renewcommand{\arraystretch}{1.05}
\resizebox{\columnwidth}{!}{
\begin{tabular}{lccc}
\toprule
\textbf{Dataset} & \textbf{\#Intents} & \textbf{\#Samples} & \textbf{Type} \\
\midrule
CLINC150 & 150 & 22500 & Multiple domains \\
StackOverflow & 20 & 20000 & Technical \\
Banking77 & 77 & 13083 & Banking \\
\bottomrule
\end{tabular}
}
\caption{Dataset statistics.}
\label{tab:dataset_stats}
\end{table}

\subsection{Baselines and Experimental Settings}
This work defines known intent ratio (KIR) to control the proportion of intent classes available during training:
\begin{equation}
\mathrm{KIR}=\frac{|\mathcal{Y}_{\mathrm{K}}|}{|\mathcal{Y}|},
\end{equation}
where $\mathcal{Y}$ denotes the complete intent set and $\mathcal{Y}_{\mathrm{K}}$ denotes the known intent set.

We evaluate KIR values of 0.25, 0.50, and 0.75. A larger KIR means that more intent classes are used for training and validation. The remaining classes are treated as OOS intents and used only for testing.

Moreover, the following representative methods are selected as baselines:

\textbf{MSP}~\citep{hendrycks2017baseline} uses the maximum softmax probability as the confidence score for detecting misclassified or out-of-distribution samples.

\textbf{OpenMax}~\citep{bendale2016towards} recalibrates activation scores and estimates the probability that an input belongs to an unknown class.

\textbf{DOC}~\citep{shu2017doc} replaces the softmax layer with independent sigmoid classifiers and performs class-wise open-set rejection.

\textbf{DeepUnk}~\citep{lin2019deep} uses margin loss to learn discriminative intent features and applies novelty detection for unknown intent detection.

\textbf{KNNCL}~\citep{zhou2022knn} combines KNN-based contrastive representation learning with nearest-neighbor information for OOD intent detection.

\textbf{ADB}~\citep{zhang2021adaptive} learns adaptive class-specific decision boundaries for open intent detection.

\textbf{DA-ADB}~\citep{zhang2023learning} improves ADB with distance-aware representation learning and adaptive boundary learning.

\textbf{MOGB}~\citep{li2025multi} models each known intent with adaptive granular-ball representations and constructs multiple local decision boundaries.

Moreover, the Gate uses the 22M-parameter \texttt{all-MiniLM-L6-v2}, while the Router and Expert use \texttt{SmolLM-135M}. For the Gate, only the last two Transformer blocks and the residual projection are learnable, with learning rates of $2\times10^{-5}$ and $2\times10^{-4}$, respectively. The intra-class and inter-class weights in Eq.~(\ref{eq:rep_loss}) are both set to 0.1, with temperature 0.07 and margin 0.20. For the Router and Expert, the LoRA ranks are 32 and 16, respectively, and the learning rate is $2\times10^{-4}$. All main experiments use only known intent samples for training and validation. All main results are averaged over three random seeds, i.e., 13, 42 and 87.


\subsection{Evaluation Metrics}

To evaluate and compare the performances of all methods, known F1 score, OOS F1 score and Acc are selected as metrics.

\textbf{known F1 score}~\citep{wang2021texsmart,zawbaa2024improved} is calculated by averaging F1 scores over all known intent classes, as follows:
\begin{equation}
\mathrm{known\ F1}=\frac{1}{|\mathcal{Y}_{\mathrm{K}}|}\sum_{y\in\mathcal{Y}_{\mathrm{K}}}\frac{2P_yR_y}{P_y+R_y},
\end{equation}
where $\mathcal{Y}_{\mathrm{K}}$ denotes the known intent label set, and $P_y$ and $R_y$ denote the precision and recall of class $y$ respectively. The precision is computed over all open-set test predictions. OOS samples predicted as known are counted as false positives for the predicted known classes.

\textbf{OOS F1 score}~\citep{wang2021texsmart,zawbaa2024improved} is calculated for evaluating OOS detection performance, as follows:
\begin{equation}
\mathrm{OOS\ F1}=\frac{2P_{\mathrm{OOS}}R_{\mathrm{OOS}}}{P_{\mathrm{OOS}}+R_{\mathrm{OOS}}},
\end{equation}
where
\begin{equation}
\begin{aligned}
P_{\mathrm{OOS}}&=\frac{\mathrm{TP}_{\mathrm{OOS}}}{\mathrm{TP}_{\mathrm{OOS}}+\mathrm{FP}_{\mathrm{OOS}}},\\
R_{\mathrm{OOS}}&=\frac{\mathrm{TP}_{\mathrm{OOS}}}{\mathrm{TP}_{\mathrm{OOS}}+\mathrm{FN}_{\mathrm{OOS}}}.
\end{aligned}
\end{equation}
$\mathrm{TP}_{\mathrm{OOS}}$, $\mathrm{FP}_{\mathrm{OOS}}$ and $\mathrm{FN}_{\mathrm{OOS}}$ denote true positives, false positives and false negatives w.r.t. the OOS class, respectively.

\textbf{Acc} is the overall accuracy over all test samples, as follows:
\begin{equation}
\mathrm{Acc}=\frac{1}{N}\sum_{i=1}^{N}\mathbb{I}(\hat{y}_i=y_i),
\end{equation}
where $N$ is the size of the test set, $y_i$ and $\hat{y}_i$ are the ground truth and predicted labels of the $i$-th sample, respectively, and $\mathbb{I}(\cdot)$ is the indicator function.

\subsection{Performance on OOS Intent Detection}

The performance on OOS intent detection is presented in Table~\ref{tab:main_results_all}. The results show that the proposed method achieves the best OOS F1 scores across all nine settings. Compared with the strongest baseline in each setting, the proposed method achieves improvements of 1.04--9.55 percentage points on OOS F1.

As shown in Table~\ref{tab:main_results_all}, the OOS F1 scores of the proposed method decrease as KIR increases on all three datasets. This indicates that OOS detection becomes more difficult when more intents are treated as known. But the proposed method consistently achieves the best OOS F1 under different KIRs and datasets. A possible reason is that the learnable MiniLM produces more compact representations for known intents. Thus, the learned boundary can better distinguish supported inputs from OOS inputs.

Moreover, the proposed method does not consistently achieve the best Known F1, although it achieves the best accuracy in most settings. This shows that the main advantage of our method lies in OOS rejection rather than known intent classification. But the cascade workflow separates these two processes. Thus, the Router--Expert component can be further improved or replaced without changing the MiniLM-based OOS detection Gate.


\begin{table*}[!t]
\centering
\fontsize{7.5}{6.9}\selectfont
\setlength{\tabcolsep}{2pt}
\renewcommand{\arraystretch}{1}

\begin{tabular*}{\textwidth}{
@{\extracolsep{\fill}}
c l
ccc
ccc
ccc
@{}
}
\toprule
&
& \multicolumn{3}{c}{\textbf{CLINC150}}
& \multicolumn{3}{c}{\textbf{StackOverflow}}
& \multicolumn{3}{c}{\textbf{Banking77}} \\
\cmidrule(lr){3-5}
\cmidrule(lr){6-8}
\cmidrule(lr){9-11}

\textbf{KIR}
& \textbf{Method}
& \textbf{Known F1} & \textbf{OOS F1} & \textbf{Acc}
& \textbf{Known F1} & \textbf{OOS F1} & \textbf{Acc}
& \textbf{Known F1} & \textbf{OOS F1} & \textbf{Acc} \\
\midrule

\multirow{9}{*}{0.25}
& MSP
& 58.43$\pm$1.30 & 73.73$\pm$0.64 & 66.25$\pm$0.55
& 43.09$\pm$1.41 & 22.92$\pm$2.71 & 32.74$\pm$1.19
& 56.05$\pm$3.06 & 63.97$\pm$10.72 & 59.16$\pm$8.87 \\

& OpenMax
& 68.55$\pm$1.40 & 86.52$\pm$0.74 & 80.49$\pm$0.88
& 46.49$\pm$0.46 & 39.47$\pm$1.95 & 41.48$\pm$1.27
& 60.10$\pm$1.03 & 75.18$\pm$7.87 & 68.84$\pm$7.63 \\

& DOC
& 7.39$\pm$4.00 & 89.35$\pm$0.28 & 80.94$\pm$0.56
& 52.87$\pm$1.31 & 60.07$\pm$2.26 & 55.22$\pm$1.75
& 63.63$\pm$1.95 & 78.53$\pm$5.88 & 72.13$\pm$5.94 \\

& DeepUnk
& 70.07$\pm$2.09 & 87.25$\pm$1.36 & 81.20$\pm$1.82
& 45.37$\pm$0.86 & 44.04$\pm$6.30 & 43.68$\pm$3.91
& 51.17$\pm$1.85 & 61.47$\pm$3.29 & 55.18$\pm$2.35 \\

& KNNCL
& 55.11$\pm$1.17 & 73.44$\pm$2.19 & 64.95$\pm$2.32
& 39.12$\pm$0.23 & 8.93$\pm$3.25 & 25.09$\pm$1.49
& 40.45$\pm$3.18 & 39.00$\pm$8.98 & 38.15$\pm$6.60 \\

& ADB
& 72.97$\pm$1.49 & 89.70$\pm$0.76 & 84.54$\pm$1.10
& 78.29$\pm$1.02 & 92.19$\pm$0.76 & 88.30$\pm$1.03
& 62.93$\pm$4.83 & 78.90$\pm$3.27 & 72.00$\pm$4.03 \\

& DA-ADB
& 77.15$\pm$1.10 & 92.87$\pm$0.27 & 88.82$\pm$0.44
& \textbf{83.02$\pm$1.96} & 94.46$\pm$1.02 & 91.52$\pm$1.50
& 65.61$\pm$3.79 & 84.46$\pm$1.33 & 77.81$\pm$2.10 \\

& MOGB
& 62.79$\pm$1.04 & 93.54$\pm$0.14 & 89.03$\pm$0.24
& 66.99$\pm$6.64 & 90.82$\pm$1.40 & 85.58$\pm$2.30
& 60.80$\pm$3.98 & 90.73$\pm$0.72 & 85.01$\pm$1.21 \\

& \textbf{Ours}
& \textbf{78.41$\pm$1.63}
& \textbf{95.19$\pm$0.48}
& \textbf{91.18$\pm$0.65}
& 80.27$\pm$0.37
& \textbf{95.50$\pm$0.96}
& \textbf{91.76$\pm$0.87}
& \textbf{76.49$\pm$2.19}
& \textbf{91.78$\pm$0.58}
& \textbf{87.36$\pm$0.98} \\

\midrule

\multirow{9}{*}{0.50}
& MSP
& 74.56$\pm$1.95 & 69.06$\pm$5.73 & 70.24$\pm$3.80
& 69.20$\pm$2.12 & 41.05$\pm$6.92 & 58.27$\pm$2.56
& 75.38$\pm$0.64 & 65.60$\pm$3.67 & 69.98$\pm$2.10 \\

& OpenMax
& 80.97$\pm$0.36 & 83.16$\pm$0.90 & 81.08$\pm$0.75
& 71.39$\pm$1.65 & 49.74$\pm$4.50 & 61.93$\pm$1.77
& 74.93$\pm$1.74 & 72.56$\pm$1.13 & 73.31$\pm$0.57 \\

& DOC
& \textbf{84.34$\pm$0.34} & 88.70$\pm$0.69 & 86.28$\pm$0.57
& 76.27$\pm$0.98 & 67.96$\pm$6.21 & 71.79$\pm$3.52
& 5.25$\pm$4.17 & 67.86$\pm$0.51 & 52.10$\pm$1.17 \\

& DeepUnk
& 82.93$\pm$0.10 & 87.78$\pm$0.40 & 85.08$\pm$0.36
& 69.00$\pm$0.95 & 39.10$\pm$1.69 & 56.78$\pm$0.70
& 65.34$\pm$2.71 & 49.83$\pm$3.74 & 57.87$\pm$2.44 \\

& KNNCL
& 81.94$\pm$0.57 & 87.21$\pm$0.32 & 84.36$\pm$0.44
& 61.83$\pm$1.20 & 7.82$\pm$0.91 & 45.11$\pm$0.75
& 53.32$\pm$2.21 & 35.93$\pm$2.22 & 46.92$\pm$0.82 \\

& ADB
& 83.34$\pm$0.35 & 87.78$\pm$0.86 & 85.24$\pm$0.76
& 84.77$\pm$0.66 & 86.42$\pm$0.28 & 85.40$\pm$0.44
& 75.58$\pm$1.75 & 69.78$\pm$2.40 & 71.54$\pm$2.14 \\

& DA-ADB
& 82.74$\pm$0.90 & 89.84$\pm$0.29 & 86.98$\pm$0.40
& \textbf{85.97$\pm$0.82}
& 87.47$\pm$0.69
& \textbf{86.54$\pm$0.80}
& 75.82$\pm$1.57 & 78.46$\pm$1.57 & 76.69$\pm$1.68 \\

& MOGB
& 60.75$\pm$0.33 & 84.63$\pm$0.15 & 78.09$\pm$0.24
& 67.51$\pm$4.71 & 79.62$\pm$1.54 & 75.01$\pm$2.62
& 63.37$\pm$1.21 & 78.47$\pm$0.56 & 73.05$\pm$0.79 \\

& \textbf{Ours}
& 82.98$\pm$0.79
& \textbf{92.13$\pm$0.81}
& \textbf{88.23$\pm$0.87}
& 83.13$\pm$1.50
& \textbf{89.23$\pm$2.04}
& 86.33$\pm$1.59
& \textbf{81.20$\pm$0.89}
& \textbf{86.39$\pm$0.99}
& \textbf{83.77$\pm$1.11} \\

\midrule

\multirow{9}{*}{0.75}
& MSP
& 86.39$\pm$0.22 & 77.09$\pm$0.16 & 82.28$\pm$0.18
& 80.57$\pm$0.33 & 37.38$\pm$3.00 & 72.51$\pm$0.53
& 84.24$\pm$0.63 & 58.04$\pm$2.38 & 78.03$\pm$0.54 \\

& OpenMax
& 84.37$\pm$0.17 & 81.86$\pm$0.12 & 83.73$\pm$0.17
& 83.17$\pm$0.46 & 54.47$\pm$6.23 & 76.63$\pm$1.28
& 80.73$\pm$1.43 & 61.15$\pm$0.80 & 76.17$\pm$0.95 \\

& DOC
& 1.51$\pm$0.60 & 58.41$\pm$0.09 & 41.54$\pm$0.20
& 85.60$\pm$0.11 & 70.66$\pm$3.29 & 81.62$\pm$0.87
& 3.04$\pm$1.89 & 39.96$\pm$0.24 & 25.98$\pm$0.82 \\

& DeepUnk
& 82.63$\pm$0.59 & 76.83$\pm$2.69 & 79.84$\pm$1.51
& 78.60$\pm$1.16 & 35.46$\pm$5.00 & 70.12$\pm$1.54
& 73.23$\pm$0.15 & 38.08$\pm$5.77 & 66.97$\pm$1.12 \\

& KNNCL
& 73.63$\pm$1.26 & 66.06$\pm$2.08 & 70.54$\pm$1.33
& 74.15$\pm$0.90 & 9.21$\pm$3.81 & 62.89$\pm$0.81
& 60.42$\pm$1.43 & 34.52$\pm$6.63 & 57.34$\pm$0.82 \\

& ADB
& \textbf{87.13$\pm$0.09} & 82.98$\pm$0.96 & 85.08$\pm$0.48
& 86.47$\pm$0.36 & 73.84$\pm$0.83 & 82.52$\pm$0.30
& 83.20$\pm$0.42 & 59.19$\pm$2.61 & 77.14$\pm$0.52 \\

& DA-ADB
& 84.58$\pm$0.29 & 84.04$\pm$0.73 & 84.53$\pm$0.51
& \textbf{87.10$\pm$0.32}
& 74.69$\pm$0.90
& \textbf{83.22$\pm$0.34}
& 79.48$\pm$0.34 & 63.69$\pm$0.55 & 74.91$\pm$0.15 \\

& MOGB
& 60.15$\pm$0.44 & 71.40$\pm$0.20 & 67.20$\pm$0.28
& 68.54$\pm$0.79 & 57.60$\pm$0.35 & 64.36$\pm$0.29
& 61.75$\pm$2.12 & 54.08$\pm$0.97 & 58.85$\pm$1.55 \\

& \textbf{Ours}
& 84.71$\pm$0.64
& \textbf{86.80$\pm$1.55}
& \textbf{85.58$\pm$1.07}
& 85.04$\pm$1.06
& \textbf{75.94$\pm$1.93}
& 82.28$\pm$0.96
& \textbf{84.53$\pm$0.77}
& \textbf{73.24$\pm$3.42}
& \textbf{81.35$\pm$1.18} \\

\bottomrule
\end{tabular*}

\caption{Performance comparison under different known intent ratios (KIRs). All results are reported as mean$\pm$standard deviation over three random seeds. The best mean result in each setting is highlighted in bold.}
\label{tab:main_results_all}
\end{table*}

\subsection{Interpretability Analysis for Gate Stage}
 
\begin{figure}[ht]
    \centering
    \includegraphics[width=\columnwidth]
    {figures/paper_style_local_boundary_geometry_3d.png}
    \caption{Local boundary geometry of Frozen and learnable MiniLM on StackOverflow at KIR$=0.50$. The surface denotes $s(\bar{e})=1$.}
    \label{fig:local_boundary_geometry}
\end{figure}

\begin{figure}[t]
    \centering
    \includegraphics[width=\columnwidth]{figures/paper_style_score_distribution.png}
    \caption{Gate score distributions on StackOverflow at KIR$=0.50$. The dashed line denotes $s(\bar{e})=1$.}
    \label{fig:gate_score_distribution}
\end{figure}

Figure~\ref{fig:local_boundary_geometry} shows the local boundary on StackOverflow with KIR$=0.50$. With Frozen MiniLM, the known and OOS samples have more overlap around the decision boundary. After training, the known representations become more compact, while more OOS samples are distributed outside the boundary. This indicates that training MiniLM with known intents can learn a more discriminative representation and provide clearer regions for OOS rejection.

Moreover, Figure~\ref{fig:gate_score_distribution} presents the corresponding Gate score distributions. After training MiniLM, the known and OOS score distributions become more separable, and more OOS samples fall into the rejection region. Since only known intent samples are used for training, the improvement shows that a better representation of known intents can also benefit the detection of unseen OOS samples. Thus, the learnable MiniLM provides a more suitable representation space for the subsequent boundary learning.

\subsection{MiniLM Selection for Gate Stage}
The ablation studies for MiniLM selection of the Gate stage are conducted, as shown in Table~\ref{tab:ablation_all_kir}. \textbf{Frozen MiniLM} means that the MiniLM encoder in the Gate is fixed without representation adaptation. \textbf{Without Gate} means to cancel the gate stage in our full pipeline, and instead conduct OOS rejection in the router module. \textbf{Cascade-MiniLM} and \textbf{Cascade-SmolLM} use MiniLM and SmolLM for the cascade modules, respectively.

\begin{table}[t]
\centering
\footnotesize
\setlength{\tabcolsep}{4.2pt}
\renewcommand{\arraystretch}{0.96}

\begin{tabular*}{\columnwidth}{
@{\extracolsep{\fill}}
lccc
@{}
}
\toprule
\multirow{2}{*}{\textbf{Variant}}
& \multicolumn{3}{c}{\textbf{KIR}} \\
\cmidrule(lr){2-4}
& \textbf{0.25}
& \textbf{0.50}
& \textbf{0.75} \\
\midrule

% ---------- CLINC150 ----------
& \multicolumn{3}{c}{\textbf{CLINC150}} \\
\cmidrule(lr){2-4}

\textbf{Ours}
& \textbf{95.19 / 91.18}
& \textbf{92.13 / 88.23}
& \textbf{86.80 / 85.58} \\

Frozen
& 90.52 / 84.39
& 84.40 / 80.18
& 72.72 / 76.08 \\

w/o Gate
& 61.96 / 54.15
& 66.33 / 66.13
& 60.93 / 71.52 \\

Casc.-MiniLM
& 90.52 / 85.58
& 84.40 / 82.25
& 72.72 / 78.53 \\

Casc.-SmolLM
& 77.25 / 68.34
& 64.36 / 64.08
& 60.85 / 70.42 \\

\midrule

% ---------- StackOverflow ----------
& \multicolumn{3}{c}{\textbf{StackOverflow}} \\
\cmidrule(lr){2-4}

\textbf{Ours}
& \textbf{95.50 / 91.76}
& \textbf{89.23 / 86.33}
& \textbf{75.94 / 82.28} \\

Frozen
& 91.60 / 86.08
& 76.83 / 75.51
& 63.51 / 76.63 \\

w/o Gate
& 56.22 / 51.41
& 58.10 / 64.50
& 54.01 / 75.00 \\

Casc.-MiniLM
& 91.60 / 86.69
& 76.83 / 76.78
& 63.51 / 78.12 \\

Casc.-SmolLM
& 21.41 / 29.38
& 18.14 / 45.44
& 14.93 / 62.80 \\

\midrule

% ---------- Banking77 ----------
& \multicolumn{3}{c}{\textbf{Banking77}} \\
\cmidrule(lr){2-4}

\textbf{Ours}
& \textbf{91.78 / 87.36}
& \textbf{86.39 / 83.77}
& \textbf{73.24 / 81.35} \\

Frozen
& 80.98 / 70.74
& 70.48 / 61.80
& 65.72 / 61.13 \\

w/o Gate
& 73.64 / 62.04
& 74.34 / 66.05
& 73.03 / 67.58 \\

Casc.-MiniLM
& 80.98 / 71.08
& 70.48 / 62.63
& 65.72 / 62.45 \\

Casc.-SmolLM
& 40.74 / 32.39
& 43.52 / 40.43
& 48.60 / 48.99 \\

\bottomrule
\end{tabular*}

\caption{
Ablation results under different known intent ratios (KIRs).
Each entry reports OOS F1 / Acc (\%).
The best result in each setting is highlighted in bold.
}
\label{tab:ablation_all_kir}
\end{table}

As shown in Table~\ref{tab:ablation_all_kir}, the full method consistently outperforms the ablation variants on OOS F1. The performance drops after removing the Gate, which shows that an explicit OOS rejection stage is important for the cascade workflow. Moreover, Frozen MiniLM performs worse than the learnable MiniLM in all settings, showing the benefit of adapting the representation with known intent samples. The results of Casc.-MiniLM and Casc.-SmolLM further show that simply using the same backbone for all stages cannot obtain the same performance. Thus, using a lightweight learnable MiniLM for OOS rejection and SmolLM for the subsequent known intent prediction is more suitable for the proposed cascade workflow.

Moreover, we first use the 135M-parameter \texttt{SmolLM-135M} for all stages actually. But after that, we replace it with the 22M-parameter \texttt{all-MiniLM-L6-v2} for gate stage, and the performance is improved. That means, it is not essential to use models with large parameters, which refers to our insight on LLMs in section Related Work.

Then, we analyze the effect of the centroid number $K$. Figure~\ref{fig:k_ablation} reports the OOS F1 with $K$ varying from 1 to 5 under KIR$=0.50$.

As shown in Figure~\ref{fig:k_ablation}, the effect of $K$ varies across different datasets. The performance on CLINC150 gradually decreases as $K$ increases, while Banking77 is relatively stable. Moreover, StackOverflow shows a more obvious performance degradation with larger $K$. This indicates that increasing the centroid number does not consistently improve OOS detection. Although multiple centroids can model more local regions of known intents, they may also enlarge the overall acceptance region and introduce more OOS samples. Thus, the appropriate boundary complexity depends on the data distribution, rather than simply using more centroids.

\begin{figure}[t]
    \centering
    \includegraphics[width=\columnwidth]{figures/k_ablation_compact.png}
    \caption{OOS F1 with different centroid numbers $K$ under KIR$=0.50$.}
    \label{fig:k_ablation}
\end{figure}

\section{Conclusion}
This work proposes a multi-cluster boundary learning method for OOS intent detection via a cascade MiniLM workflow. That workflow separates the OOS intent detection from the multi-class classification of known intents, and views it as a one-class classification problem. Then, multi-cluster boundary learning is conducted to learn the boundary of MiniLM embeddings from the input user utterances. That boundary separates the embedding domain of OOS intents and known intents. Extensive experiments are conducted on public datasets. The results show that the proposed method achieve the stete-of-the-art OOS intent detection performance. Ablation studies also show that the MiniLM settings are reasonable at present. Moreover, the performance of the method is not always the best one on the task of known intent detection, but the cascade workflow completely decouples the two detection process for OOS and known intents. Thus, the result does not influence the aforementioned conclusion. And moreover, the lightweight models have outstanding advantages for real-world applications compared to the large models. The more advanced MiniLMs perhaps improve the detection performance by using this framework in the future.       




\section*{Limitations}
The proposed method in this work achieves the state-of-the-art performance on OOS intent detection task, but does not generalize well on the task of known intent detection. Thus our primary problem setting is OOS intent detection in the title. 

\section*{Ethics Statement}
This paper does not involve text reuse from previously published work. We do not use any AI-assisted tools for writing this paper.

\bibliography{custom}

\end{document}
